\section{Introduction}\label{sec:intro}

Steel structures have historically been conceived as passive load-resisting systems. A frame, a truss or a braced industrial skeleton is designed to a code-specified set of limit states, erected, and thereafter assumed to behave as drawn for the remainder of its service life. Within this paradigm, knowledge of the structure's actual condition is acquired discontinuously, through periodic visual inspection and occasional non-destructive evaluation, and is mediated entirely by human judgement. The structure itself contributes nothing to the assessment of its own integrity: it carries load, accumulates damage silently, and waits to be examined. This passive model has well-known limitations. Inspections are infrequent, labour-intensive and costly; they favour accessible and visible regions; and they capture the state of the structure only at isolated instants, leaving long intervals during which deterioration, fatigue crack growth, support settlement or the progressive loosening of bolted connections may develop undetected \citep{avci2021review}.

Structural Health Monitoring (SHM) emerged precisely to overcome the discontinuity of inspection. SHM is the process of implementing a damage detection strategy for aerospace, civil and mechanical infrastructure by observing a structure over time through periodically sampled response measurements, extracting damage-sensitive features and statistically analysing those features to determine the current state of system health \citep{farrar2013shm,farrar2007introduction}. The premise is physically grounded: damage is a change in the material or geometric properties of a system, including changes to boundary conditions and connectivity, that alters its stiffness, mass or energy-dissipation characteristics and therefore its measured dynamic response \citep{farrar2013shm}. The natural vibration properties of a structure depend only on its mass and stiffness, so that its natural frequencies and mode shapes are fixed by these two quantities alone \citep{chopra2017dynamics}; any stiffness loss is, in principle, observable as a measurable shift in frequencies, mode shapes and damping. The textbook precedent of period lengthening accompanying stiffness reduction confirms that monitored modal parameters genuinely track structural condition \citep{chopra2017dynamics}.

Yet most deployed SHM systems stop short of their own promise. They excel at instrumentation, acquisition, storage and transmission of data, but the interpretation of those data, the conversion of raw signals into a statement about the structure's health, is too often performed off the structure, by human experts and after the fact. Automated and near-real-time SHM methods do exist \citep{avci2021review,he2022integrated}; what remains uncommon is interpretation embedded in the structure itself, together with an explicit measure of how autonomously it acts. This reflects a structural imbalance in the field itself: feature extraction and, especially, the statistical modelling that discriminates damage states have historically received far less attention than data acquisition, even though it is here that sensor data become damage information \citep{farrar2013shm}. A sensor cannot measure damage directly; without intelligent feature extraction and statistical classification, the structure remains mute \citep{farrar2013shm}. The result is a generation of monitored structures that are richly instrumented but cognitively inert: they sense, but they do not understand.

A convergence of enabling technologies now makes it feasible to close this gap. Embedded and distributed sensing, using dense arrays of strain gauges, triaxial accelerometers, temperature sensors, displacement transducers and inclinometers fused across heterogeneous channels \citep{farrar2013shm}, provides continuous, multi-physics observation. The finite element method supplies a rigorous physical baseline: assembling element stiffness matrices into a global model and solving the free-vibration eigenproblem yields the natural frequencies and mode shapes against which measured behaviour can be compared \citep{cook2002fea}. The FE model is simulation, not reality, and its disagreement with experiment is itself diagnostic \citep{cook2002fea}, a property that turns model--data correlation into a continuous deviation detector. Machine learning supplies the interpretive layer: supervised methods such as Random Forest, Support Vector Machines and Artificial Neural Networks classify damage type and severity, while unsupervised methods such as autoencoders and recurrent LSTM networks learn healthy behaviour and flag novelty \citep{naser2023ml,brunton2019data,wangcha2022unsupervised}. Data-driven dimensionality reduction (SVD/PCA), equation-free modal extraction (DMD) and physics-informed learning provide the algorithmic machinery to fuse model and measurement into a single predictive object, the digital twin \citep{brunton2019data,chiachio2022structural}. Embedded within an Industry 4.0 ecosystem of continuous acquisition, edge processing and connectivity, these elements together permit a permanent, real-time link between a physical structure and its digital representation \citep{li2025edge,zou2026deep}.

Building on this convergence, we introduce the concept of the \textbf{Self-Aware Steel Structure}: a structure able to perceive, process, interpret and continuously report its own mechanical state. We name the resulting paradigm \textbf{Structural Cognition}. Under it, a steel structure ceases to be merely a load-resisting element and becomes an intelligent system that interprets its own behaviour, recognises changes in operational condition, identifies stiffness loss and bolted-connection loosening, detects progressive damage, and autonomously assesses its integrity, advancing along the damage-state hierarchy from existence and location toward type, severity and prognosis \citep{farrar2013shm}.

The research gap addressed by this paper is conceptual and architectural rather than purely instrumental. Digital twins, finite-element models and machine learning have already been combined for structural monitoring, including structural digital-twin formulations \citep{chiachio2022structural}, deep-learning- and digital-twin-driven SHM \citep{zou2026deep}, and autonomous, distributed monitoring architectures \citep{li2025edge,mariniello2025enhancing}. What these works do not provide is an explicit, operational scale of \textit{how autonomously} a structure interprets and acts on its own state, together with a steel-specific architecture and protocol that operationalise it. The objective of this paper is to supply that missing layer, whose architecture is summarised in Fig.~\ref{fig:arch}. This is a conceptual framework and review paper: it proposes the Structural Cognition paradigm, an operational six-level taxonomy (L0--L5) of structural self-awareness, a reference architecture for Self-Aware Steel Structures, and a detailed experimental specification centred on an instrumented three-dimensional braced frame of square hollow sections excited by a controlled rotating-unbalance system. The paper does not report empirical measurements; the experimental campaign is described as a planned validation pathway. Its contributions are threefold: (i) the formal articulation of Structural Cognition together with an operational L0--L5 maturity scale that moves SHM from data acquisition toward autonomous interpretation; (ii) a reference architecture integrating multi-sensor fusion, a continuously updated digital twin and a supervised/unsupervised learning layer grounded in established theory, whose five foundational pillars are listed in Table~\ref{tab:pillars} \citep{farrar2013shm,brunton2019data,chopra2017dynamics,naser2023ml,cook2002fea}; and (iii) a structured, auditable six-step methodology and experimental specification for realising and validating a Self-Aware Steel Structure.

The remainder of the paper is organised as follows. Section~\ref{sec:background} reviews the theoretical foundations of SHM, structural dynamics, finite element modelling, digital twins and machine learning that underpin Structural Cognition. Section~\ref{sec:paradigm} defines the Self-Aware Steel Structure and the Structural Cognition paradigm conceptually. Section~\ref{sec:architecture} presents the proposed reference architecture. Section~\ref{sec:methodology} details the experimental methodology and instrumented testbed. Section~\ref{sec:ml} describes the digital-twin implementation and the learning layer for autonomous integrity assessment. Section~\ref{sec:discussion} discusses implications, limitations and open challenges, and Section~\ref{sec:conclusions} concludes.
